% GENERATED FILE. Edit canonical lessons, then run npm run generate:books.
% canonical-source-hash: fnv1a64:d7715072b975cbc3
% canonical-lessons: ES-C33-verde-amarillo

\chapter{Green and Yellow}
\label{ch:spanish-green-and-yellow}

This chapter is generated from the canonical micro-lessons used by Language Ladder.
Each section stays independently resumable and preserves the authored prerequisite order.

\section[verde, amarillo]{verde, amarillo — the expected inheritance, and a real surprise}

\label{lesson:ES-C33-verde-amarillo}



\begin{quote}
\textbf{Warm-up.} {[}PAUSE 2s{]} One of these colors is exactly what you'd expect from Latin. The other one will genuinely surprise you — it has nothing to do with any Latin word for \textquotedblleft{}yellow\textquotedblright{} at all.
\end{quote}

\begin{cousinweb}
\textbf{verde} (\textquotedblleft{}\textbf{green}\textquotedblright{}) $\leftarrow$ Latin \textbf{viridis} (worn down via \textbf{virdis}), from the verb \textbf{virēre}, \textquotedblleft{}\textbf{to be green, to flourish, to be vigorous}.\textquotedblright{} This is a cousin of English \textbf{verdant} and \textbf{verdure} — though \textbf{not} of English \textbf{green} itself, which comes from a completely separate Germanic root. A clean, expected Latin inheritance.
\end{cousinweb}

\begin{cousinweb}
Here's the honest surprise: \textbf{amarillo} (\textquotedblleft{}\textbf{yellow}\textquotedblright{}) does \textbf{not} trace to any of Latin's actual words for yellow (\emph{flāvus}, \emph{gilvus}, \emph{lūteus}). Instead, it comes from Late Latin \textbf{amarellus}, \textquotedblleft{}\textbf{a little bitter}\textquotedblright{} — a diminutive of \textbf{amarus}, \textquotedblleft{}\textbf{bitter}.\textquotedblright{} The connection runs through an old association between the color yellow and \textbf{bitterness} — bile, jaundice, and many bitter-tasting plants all share a yellowish cast. Over centuries, \textquotedblleft{}a little bitter\textquotedblright{} softened into simply \textquotedblleft{}yellowish,\textquotedblright{} and by Spanish's Golden Age the word had shed its bitter meaning entirely, leaving just the color. \emph{Amarillo} is a genuine cousin of Spanish's own \textbf{amargo} (\textquotedblleft{}bitter\textquotedblright{}) — both descend from the same \emph{amarus}.
\end{cousinweb}

\subsection*{Guided Practice}
{[}PAUSE 1s{]}

\begin{itemize}
\raggedright
  \item {[}YOU SAY: \textquotedblleft{}verde\textquotedblright{} — green, from \textquotedblleft{}to flourish\textquotedblright{}{]}
  \item {[}YOU SAY: \textquotedblleft{}amarillo\textquotedblright{} — yellow, literally \textquotedblleft{}a little bitter\textquotedblright{}{]}
  \item {[}YOU SAY: the cousins — \textquotedblleft{}amarillo, amargo\textquotedblright{} — yellow and bitter, same Latin root{]}
\end{itemize}

\begin{tcolorbox}[breakable,colback=teal!4,colframe=teal!35!black,title={Wrap-up recall}]
{[}PAUSE 3s{]} What Latin verb gives \textbf{verde}, and what does it mean? (\textbf{Virēre}, \textquotedblleft{}to be green/flourish\textquotedblright{} — cousin of English \emph{verdant}, \emph{verdure}, not of English \emph{green} itself.) Does \textbf{amarillo} come from a Latin \textquotedblleft{}yellow\textquotedblright{} word? (\textbf{No} — from \emph{amarellus}, \textquotedblleft{}a little bitter,\textquotedblright{} a diminutive of \emph{amarus}, \textquotedblleft{}bitter.\textquotedblright{}) What everyday Spanish word is \textbf{amarillo}'s cousin? (\textbf{Amargo}, \textquotedblleft{}bitter.\textquotedblright{})
\end{tcolorbox}
